% ============================================================
\chapter{Contexto y estado del arte}
\label{cap:estado}
% ============================================================

La detección automática de vulnerabilidades en infraestructuras \textit{Kubernetes}
se encuentra en la intersección de cuatro áreas de investigación activa: el análisis
de seguridad en infraestructura como código (\textit{Infrastructure as Code}, \acs{IaC}),
el aprendizaje automático aplicado a la ciberseguridad, las redes neuronales de grafos
para razonamiento estructural, y la construcción de grafos de ataque para la
identificación de rutas de explotación. A continuación se revisan los trabajos y
herramientas más relevantes en cada área, y se posiciona la contribución de este
trabajo respecto a ellos.

% -----------------------------------------------------------
\section{Contexto del problema: seguridad en infraestructuras Kubernetes}
% -----------------------------------------------------------

\textit{Kubernetes} es la plataforma de orquestación de contenedores dominante en
entornos de producción empresariales, con un 96\,\% de organizaciones que la utilizan
o evalúan según la encuesta anual de 2021 de la \ac{CNCF}
\citep{CNCF2022}. A pesar de su adopción generalizada, la seguridad sigue siendo uno
de los principales retos operacionales: el informe \textit{State of Kubernetes Security}
de Red Hat 2024 \citep{RedHat2024} revela que el 89\,\% de las organizaciones sufrió
al menos un incidente de seguridad relacionado con \textit{Kubernetes} en el último
año, y que el 67\,\% retrasó despliegues a producción por preocupaciones de seguridad.
El 46\,\% reportó pérdidas de clientes o ingresos directamente atribuibles a incidentes
en entornos de contenedores. Entre los riesgos más temidos, las vulnerabilidades
representan el 33\,\% de las preocupaciones de los equipos, mientras que las
misconfiguraciones y exposiciones alcanzan el 27\,\%.

Su modelo declarativo, basado en ficheros \acs{YAML} que especifican el estado deseado del
clúster, ofrece grandes ventajas en reproducibilidad y automatización, pero también
introduce una amplia superficie de ataque: un fichero mal configurado desplegado en
producción puede conceder privilegios excesivos a un contenedor o exponer secretos
en texto plano. La flexibilidad de \textit{Kubernetes} —que permite configurar desde
el sistema de ficheros del anfitrión hasta los permisos de red a nivel de pod— implica
que el número de vectores de ataque crece proporcionalmente con la riqueza del modelo
de configuración.

La Agencia de Seguridad Nacional de Estados Unidos (\acs{NSA}) y la Agencia de
Ciberseguridad e Infraestructura (\acs{CISA}) publicaron en 2022 una guía técnica
\citep{NSA2022} que enumera las misconfiguraciones más críticas: uso de
\texttt{privileged: true}, montaje del socket de Docker, compartición del espacio de
nombres de proceso del anfitrión (\textit{hostPID}), y automontaje de tokens de cuenta
de servicio con permisos excesivos. Estos vectores son precisamente los que este trabajo
modeliza como nodos de escape en el grafo del clúster.

Los trabajos de BishopFox \citep{BishopFox2021} demostraron empíricamente que un
único pod con configuración \textit{everything-allowed} puede comprometer el nodo
anfitrión en segundos y pivotar lateralmente a través del clúster. La plataforma
educativa \textit{Kubernetes Goat} \citep{KubernetesGoat} sistematiza más de
veinte escenarios de vulnerabilidad con los que se pueden reproducir rutas de
explotación completas en un entorno controlado. Ambas fuentes son la base de las etiquetas de cadena de ataque
del presente trabajo.

El marco taxonómico de referencia para los ataques en entornos de contenedores y
\textit{Kubernetes} es \textit{MITRE ATT\&CK for Containers}
\citep{MITRE2021Containers}, incorporado en la versión~9 de la base de conocimiento
ATT\&CK. Este marco describe las tácticas y técnicas empleadas por adversarios
a nivel de orquestador (\textit{Kubernetes}) y a nivel de contenedor (\textit{Docker}),
incluyendo las que este trabajo detecta directamente: escape de contenedor hacia el
nodo anfitrión (T1611, \textit{Escape to Host}), movimiento lateral mediante el uso
de \textit{tokens} de cuentas de servicio (T1550.001, \textit{Use Alternate
Authentication Material: Application Access Token}) y abuso de cuentas válidas
(T1078, \textit{Valid Accounts}). La alineación
entre las \textit{escape flags} del presente sistema y las técnicas ATT\&CK valida que
el espacio de detección cubre los vectores de ataque documentados y más explotados en
la industria.

% -----------------------------------------------------------
\section{Análisis estático de \acs{IaC}}
% -----------------------------------------------------------

El análisis estático de manifiestos \ac{IaC} cuenta con una larga trayectoria en la
literatura. Rahman et~al.\ \citep{Rahman2022} construyeron un conjunto de datos
de referencia para manifiestos \textit{Kubernetes} extraídos de repositorios
públicos, con 2.039 manifiestos de 92 repositorios etiquetados en 11 categorías de
\textit{misconfiguration} identificadas mediante análisis cualitativo y cuantificadas
con su herramienta de análisis estático (SLI-KUBE). De
ese trabajo, el presente sistema recupera y codifica 1.906 manifiestos de GitHub y
GitLab como vectores de características binarias que indican la presencia o ausencia
de cada patrón.

Las herramientas industriales actuales más utilizadas son \textit{Checkov}
\citep{Checkov2021}, que verifica más de mil políticas de seguridad predefinidas, y
\textit{kube-linter} \citep{KubeLinter2021}, especializado en buenas prácticas de
\textit{Kubernetes}. \textit{kube-linter} opera al nivel de manifiesto individual,
mientras que \textit{Checkov} incorpora además comprobaciones basadas en grafos sobre
las dependencias entre recursos (por ejemplo, escalada de privilegios a través de un
\textit{RoleBinding}); no obstante, se trata de reglas declarativas escritas manualmente,
sin un modelo aprendido de la estructura del clúster ni predicción supervisada de
cadenas de ataque. También destaca
\textit{Kubescape} \citep{Kubescape2021}, una plataforma de seguridad de código
abierto integrada en la \ac{CNCF} que combina
análisis estático con validación de cumplimiento contra los marcos \ac{NSA}, MITRE ATT\&CK
y \acs{CIS} \textit{Benchmarks}. Al igual que Checkov y kube-linter, sin embargo, Kubescape
evalúa recursos de forma independiente y no construye ningún modelo del clúster como
estructura relacional.

La noción de \textit{Unified Misconfiguration Index} (\acs{UMI}), propuesta por
\citeauthor{Malul2024} (\citeyear{Malul2024}) en \textit{GenKubeSec}, normaliza y
unifica las etiquetas de múltiples herramientas de análisis basadas en reglas para
poder compararlas, un problema afín al de generar un \textit{risk\_score} unificado
como el de la Capa~1 del presente sistema. Su propuesta, que incorpora un modelo de
lenguaje de gran escala (\acs{LLM}) para la clasificación, es complementaria a la de
este trabajo, aunque opera en el mismo nivel de manifiesto individual.

La limitación estructural común a estas herramientas es que ninguna aprende ni predice
cadenas de ataque a partir de la estructura del clúster: incluso las comprobaciones de
grafo de \textit{Checkov} son reglas fijas sobre dependencias conocidas, no un modelo
entrenado. Estas herramientas detectan que un pod tiene \texttt{privileged: true}, pero no infieren
si ese pod encadena con el acceso de red a los recursos críticos del clúster para formar
una ruta de ataque completa. La detección de cadenas de ataque requiere, inevitablemente,
razonar sobre la estructura del clúster como un todo.

% -----------------------------------------------------------
\section{Aprendizaje automático para detección de vulnerabilidades}
% -----------------------------------------------------------

El aprendizaje automático se ha aplicado a la detección de vulnerabilidades en
código fuente y configuraciones desde hace más de una década. Los clasificadores
basados en \textit{Random Forest} \citep{Breiman2001} —que combinan múltiples árboles
de decisión mediante \textit{bagging} y selección aleatoria de características— han
demostrado resultados robustos en escenarios de alta dimensionalidad y clases
desbalanceadas, características típicas de los conjuntos de datos de seguridad; el
desequilibrio se aborda además mediante la ponderación de clases que ofrece su
implementación en \textit{scikit-learn} \citep{Pedregosa2011}. El ajuste de
hiperparámetros suele realizarse mediante búsqueda en rejilla (\textit{grid search})
con validación cruzada estratificada sobre parámetros como \texttt{n\_estimators},
\texttt{max\_depth} y \texttt{min\_samples\_leaf}; los valores concretos explorados en
este trabajo se detallan en el Capítulo~\ref{cap:diseno}.

En el dominio de la detección de intrusiones en red, el conjunto de datos UNSW-NB15
\citep{Moustafa2015} se ha convertido en un referente de evaluación para clasificadores
supervisados. El conjunto de manifiestos de este trabajo presenta un desequilibrio
de clases moderado (56,3\,\% seguros vs.\ 43,7\,\% mal configurados); el
desafío de desequilibrio más severo se concentra en las \textit{flags} de escape
individuales, algunas presentes en menos del 1\,\% de los manifiestos reales.
El desequilibrio de clases es un problema generalizado
en los conjuntos de datos de seguridad: las misconfiguraciones críticas son rarísimas en
repositorios reales, lo que obliga a adoptar estrategias explícitas de compensación.
La alternativa más común es asignar pesos de clase inversamente proporcionales a su
frecuencia (\textit{class\_weight=balanced} en \textit{scikit-learn}), que es la
estrategia adoptada en este trabajo \citep{Pedregosa2011}. Otras alternativas como
\acs{SMOTE} (\textit{Synthetic Minority Over-sampling Technique}) no se consideraron
apropiadas dado que las características son binarias, y la interpolación lineal
entre muestras binarias puede generar vectores fuera del espacio semántico válido.

Trabajos más recientes han explorado el uso de modelos de lenguaje de gran escala
(\acs{LLM}) para la detección y remediación de misconfiguraciones en
\textit{Kubernetes} \citep{Malul2024}. Si bien los \acs{LLM}s ofrecen
capacidades de razonamiento semántico
superiores, requieren recursos computacionales sustancialmente mayores, producen
salidas no deterministas y son más opacos que los clasificadores clásicos, lo que
dificulta su interpretación en auditorías de seguridad donde la trazabilidad de las
reglas de detección es un requisito normativo. Para entornos de producción donde el
operador de seguridad debe poder justificar ante una auditoría por qué un clúster
fue marcado como de alto riesgo, la interpretabilidad del \textit{Random Forest}
—mediante importancia de características y reglas de árbol legibles— representa una
ventaja operacional significativa sobre los modelos de caja negra.

La elección de \textit{Random Forest} sobre otros algoritmos clásicos para la
clasificación de manifiestos \ac{IaC} responde a tres criterios documentados en la
literatura de aprendizaje automático aplicado a la seguridad:

\begin{enumerate}
\raggedright
  \item \textbf{Robustez ante datos binarios y escasos.} Las características del
    presente trabajo son en su mayoría binarias (presencia/ausencia de una
    misconfiguration), con alta proporción de ceros. Los \textit{Random Forests}
    manejan este tipo de datos de forma nativa sin requerir normalización, a diferencia
    de las Máquinas de Vectores Soporte (\acs{SVM}) o las Redes Neuronales, sensibles
    a la escala de las características.

  \item \textbf{Gestión intrínseca del desequilibrio de clases.} El parámetro
    \texttt{class\_weight=balanced} de \textit{scikit-learn} implementa una estrategia
    de ponderación inversamente proporcional a la frecuencia de clase, lo que eleva
    efectivamente el coste de clasificar erróneamente las muestras de la clase
    minoritaria (misconfiguraciones críticas) sin modificar la función de pérdida
    subyacente \citep{Pedregosa2011}.

  \item \textbf{Generación de probabilidades calibradas.} El promedio de los votos
    del \textit{ensemble} de árboles produce probabilidades suaves en \([0,1]\) que
    pueden utilizarse directamente como características continuas (\textit{risk\_scores})
    en el grafo de la Capa~2. Algoritmos como los Árboles de Decisión individuales o
    los \textit{k}-Nearest Neighbours producen probabilidades menos calibradas en
    presencia de desequilibrio de clases.
\end{enumerate}

La Tabla~\ref{tab:ml_comparison} resume esta comparativa entre algoritmos
candidatos para la Capa~1.

\begin{table}[htbp]
\centering\small
\caption{Comparativa de algoritmos de \acs{ML} para clasificación de manifiestos \ac{IaC}}
\label{tab:ml_comparison}
\begin{tabular}{@{}p{2.7cm}cp{2.1cm}p{2cm}p{4.1cm}@{}}
\toprule
Algoritmo & Binario & Desequilibrio & Prob.\ calibrada & Observación \\
\midrule
Random Forest & \checkmark & \checkmark & Buena & Adoptado (Capa~1) \\
Árbol de Decisión & \checkmark & Parcial & Pobre & Alta varianza en grafos \\
\ac{SVM} (RBF) & \checkmark & Parcial & Requiere Platt & Entrenamiento O(n\textsuperscript{2})--O(n\textsuperscript{3}); viable pero menos eficiente \\
Red Neuronal (\acs{MLP}) & \checkmark & Con pesos & Variable & Requiere normalización \\
\textit{k}-NN & \checkmark & No & Estimada & Costoso en inferencia \\
\bottomrule
\end{tabular}
\fuentepropia
\end{table}

La aplicación de técnicas de \ac{ML} a la predicción de vulnerabilidades en software
es un enfoque consolidado que supera de forma consistente a los modelos estadísticos
clásicos \citep{Jabeen2022}. Dentro de este espacio, el \textit{Random Forest} resulta
especialmente adecuado para el problema de la Capa~1 por sus propiedades sobre
características binarias de alta dimensión con clases desbalanceadas —manejo nativo de
datos dispersos, gestión intrínseca del desequilibrio y salidas probabilísticas
utilizables como \textit{risk\_score}— que se detallan a continuación. Estas
propiedades, combinadas con sus ventajas de interpretabilidad para auditorías de
seguridad, motivaron su adopción como componente de Capa~1 de \textit{kubescan}.

% -----------------------------------------------------------
\section{Redes neuronales de grafos}
% -----------------------------------------------------------

Las redes neuronales de grafos (\acs{GNN}) han emergido como la herramienta principal para
aprendizaje en datos con estructura relacional \citep{Zhou2020, Hamilton2020}. A
diferencia de los clasificadores tabulares clásicos, las \acs{GNN}s operan directamente
sobre la topología del grafo, propagando información entre vecinos mediante operaciones
de paso de mensajes (\textit{message passing}). Formalmente, en una capa \ac{GNN} la
representación del nodo \(v\) se actualiza como:

\begin{equation}
  \mathbf{h}_v^{(l+1)} = \sigma\!\left(
    \mathbf{W}^{(l)} \cdot \text{AGGREGATE}\!\left(
      \left\{ \mathbf{h}_u^{(l)} : u \in \mathcal{N}(v) \right\}
    \right) + \mathbf{b}^{(l)}
  \right)
  \label{eq:gnn_update}
\end{equation}
\addequationtoindex{Actualización de la representación de un nodo en una capa GNN}

donde \(\mathcal{N}(v)\) es el conjunto de vecinos de \(v\), \(\mathbf{W}^{(l)}\) y
\(\mathbf{b}^{(l)}\) son parámetros entrenables de la capa \(l\), y \(\sigma\) es una
función de activación no lineal. La elección de la función AGGREGATE diferencia las
principales arquitecturas.

Las arquitecturas más relevantes para este trabajo son:

\textbf{Graph Convolutional Networks (\acs{GCN})} \citep{Kipf2017}: implementan AGGREGATE
como la media ponderada de los estados de los vecinos normalizada por sus grados.
Son eficientes, pero los pesos quedan fijados por la estructura del grafo (los grados
de los nodos) y no se aprenden, sin distinción entre tipos de arista.

\textbf{GraphSAGE} \citep{Hamilton2017}: extensión inductiva que permite aplicar el
modelo a grafos no vistos durante el entrenamiento, relevante en escenarios donde
aparecen nuevos clústeres en producción. La inducción se logra aprendiendo funciones
de agregación en lugar de \textit{embeddings} fijos.

\textbf{Graph Attention Networks (\acs{GAT})} \citep{Velickovic2018}: sustituyen la
normalización fija de \ac{GCN} por coeficientes de atención aprendidos. Para cada par
de nodos \((v, u)\) el coeficiente de atención es:

\begin{equation}
  \alpha_{vu} = \frac{
    \exp\!\left( \text{LeakyReLU}\!\left(
      \mathbf{a}^\top [\mathbf{W}\mathbf{h}_v \,\|\, \mathbf{W}\mathbf{h}_u]
    \right) \right)
  }{
    \sum_{k \in \mathcal{N}(v)} \exp\!\left( \text{LeakyReLU}\!\left(
      \mathbf{a}^\top [\mathbf{W}\mathbf{h}_v \,\|\, \mathbf{W}\mathbf{h}_k]
    \right) \right)
  }
  \label{eq:gat_attn}
\end{equation}
\addequationtoindex{Coeficiente de atención en Graph Attention Networks}

donde \(\mathbf{a}\) es un vector de parámetros de atención y \(\|\) denota
concatenación. Este diseño es especialmente apropiado para grafos de clúster
\textit{Kubernetes}, donde una arista de privilegio (\textit{privilege\_reach})
debe pesar más que una arista de proximidad de directorio (\textit{dir\_proximity}).
El modelo \textit{How Powerful are \acs{GNN}s?} de \citeauthor{Xu2019} (\citeyear{Xu2019})
establece que la capacidad discriminativa de las \acs{GNN}s está acotada por el test de
isomorfismo de Weisfeiler-Lehman, y que la suma de vecinos maximiza el poder
expresivo frente a la media o el máximo. En este trabajo se adopta un
\textit{pooling} combinado de media y máximo por su robustez ante grafos de tamaño
muy variable, aceptando explícitamente ese compromiso de expresividad.

% -----------------------------------------------------------
\section{\acs{GNN}s aplicadas a ciberseguridad}
% -----------------------------------------------------------

La aplicación de \acs{GNN}s a la detección de amenazas es un campo emergente con rápido
crecimiento. \citeauthor{Bilot2023} (\citeyear{Bilot2023}) proporcionan una
revisión exhaustiva de los enfoques basados en \ac{GNN} para la detección de
intrusiones en red, organizada en torno a enfoques basados en red (grafos de flujo) y
basados en host (grafos de procedencia). Una segunda revisión publicada
en \textit{Computers and Security} en 2024 \citep{GNNSurvey2024} amplía esta
taxonomía y documenta un crecimiento sostenido de la investigación en el área en
los últimos años.

\textit{E-GraphSAGE} \citep{Lo2022} propone representar flujos de red como grafos de
transacción para clasificación de intrusiones en \acs{IoT}, alcanzando tasas de detección
superiores a los clasificadores tabulares clásicos en el conjunto de datos UNSW-NB15. La
analogía con el presente trabajo es directa: los manifiestos \textit{Kubernetes}
son los nodos y las relaciones de privilegio y movimiento lateral son las aristas del
grafo de amenaza. Una diferencia clave es que \textit{E-GraphSAGE} modela flujos de
tráfico en tiempo real, mientras que este trabajo modela configuraciones declarativas
estáticas, lo que permite operar antes del despliegue.

\citeauthor{Li2022} (\citeyear{Li2022}) construyen grafos de conocimiento de técnicas
de ataque a partir de informes de inteligencia de amenazas (\acs{CTI}) y demuestran que el
razonamiento estructural sobre relaciones entre técnicas mejora la predicción de
campañas de ataque. Este enfoque comparte con el presente trabajo la intuición central
de que el grafo de relaciones entre recursos es más informativo que la suma de sus
nodos individuales. Aunque opera sobre conocimiento de ataques pasados (retrospectivo)
y no sobre configuraciones declarativas (preventivo), la arquitectura basada en grafos
de conocimiento inspiró el diseño de los tipos de arista semánticamente diferenciados
del presente sistema.

Una limitación reconocida en la literatura es la escasez de conjuntos de datos de
grafos etiquetados de alta calidad para ciberseguridad. Tanto \citeauthor{Bilot2023}
(\citeyear{Bilot2023}) como \citeauthor{GNNSurvey2024} (\citeyear{GNNSurvey2024})
señalan la escasez de conjuntos de datos etiquetados y de benchmarks estandarizados
entre los principales obstáculos para el avance del área, lo que refuerza la contribución del presente trabajo al
generar un conjunto de datos original de 599 grafos de clúster (1.154 tras
aumentación de datos) con etiquetas de cadena de ataque.

% -----------------------------------------------------------
\section{Detección de cadenas de ataque y grafos de ataque en Kubernetes}
% -----------------------------------------------------------

La noción de cadena de ataque (\textit{kill chain}) fue formalizada originalmente
por \citeauthor{Hutchins2011} (\citeyear{Hutchins2011}) para describir las fases
secuenciales de una intrusión avanzada: reconocimiento, weaponización, entrega,
explotación, instalación, comando y control, y acción. En el contexto de
\textit{Kubernetes}, esta cadena se concreta de forma más compacta: requiere como
mínimo (1) un pod con capacidad de escape al anfitrión (\textit{TRUE\_HOST\_PID},
\textit{privileged: true}, \textit{hostPath} montado, etc.) y (2) un mecanismo de
movimiento lateral hacia otros recursos del clúster (token de cuenta de servicio con
permisos excesivos, montaje de secretos en el manifiesto, etc.).

En los últimos años han surgido dos herramientas de código abierto que abordan la
detección de rutas de ataque en \textit{Kubernetes} mediante grafos:
\textit{KubeHound} \citep{KubeHound2023} e \textit{IceKube} \citep{IceKube2023}.

\textbf{KubeHound} (Datadog, 2023) construye un grafo de ataque del clúster a partir
de la consulta directa al servidor \ac{API} de \textit{Kubernetes}, almacena el resultado
en JanusGraph y permite explorar las rutas de ataque más cortas mediante el lenguaje
de consulta Gremlin. El proyecto fue motivado por la misma intuición que subyace al
presente trabajo, formulada por John Lambert y citada en su documentación:
\textit{«los defensores piensan en listas, los atacantes piensan en grafos»}
\citep{KubeHound2023}. KubeHound cubre un
amplio conjunto de ataques y ha sido adoptado por equipos de seguridad en producción.

\textbf{IceKube} \citep{IceKube2023} sigue un enfoque análogo, basado
en Neo4j, y permite definir rutas de ataque mediante consultas Cypher. Modela un
amplio conjunto de técnicas de ataque como relaciones en el grafo y puede encontrar
rutas desde cualquier identidad de bajo privilegio hasta \textit{cluster-admin}.

La diferencia fundamental entre estas herramientas y \textit{kubescan} es de
paradigma: KubeHound e IceKube operan sobre un \textbf{clúster activo}, requiriendo
acceso con privilegios elevados a \texttt{kubectl}. \textit{kubescan}, en cambio,
opera sobre \textbf{manifiestos \acs{YAML} estáticos} antes del despliegue, integrándose
en el pipeline \acs{CICD} como un control de seguridad \textit{shift-left}. Además, ambas
herramientas utilizan búsqueda de rutas mediante consultas ad-hoc en bases de datos
de grafos, sin componente de aprendizaje automático supervisado; \textit{kubescan}
es el primer sistema, según la revisión bibliográfica realizada en este trabajo (Sección~\ref{sec:sintesis}), que combina análisis estático
de manifiestos con \acs{GNN}s para la predicción de cadenas de ataque como tarea de
clasificación supervisada.

% -----------------------------------------------------------
\section{Conclusiones del estado del arte}
\label{sec:sintesis}
% -----------------------------------------------------------

El análisis de la literatura permite identificar tres carencias fundamentales que
motivan el presente trabajo. En primer lugar, las herramientas de análisis estático
existentes (Checkov, kube-linter, Kubescape) no aprenden ni predicen cadenas de ataque
a partir de la estructura del clúster; aunque Checkov incorpora comprobaciones de grafo
basadas en reglas declarativas sobre dependencias entre recursos, ninguna emplea un
modelo supervisado que razone sobre las relaciones que habilitan dichas cadenas. En segundo lugar, los conjuntos de datos de referencia para
seguridad de \textit{Kubernetes} son escasos, de pequeño tamaño y desequilibrados
hacia las clases de riesgo bajo, lo que dificulta el entrenamiento de clasificadores
de clúster: el análisis del corpus de manifiestos reales recopilado en este trabajo
(Capítulo~\ref{cap:diseno}) confirma que \texttt{TRUE\_HOST\_PID} —la flag de escape
más peligrosa— aparece en sólo el 0,3\,\% de las filas antes de incorporar fuentes
especializadas. En tercer lugar, las herramientas de grafo de
ataque más avanzadas (KubeHound, IceKube) requieren un clúster activo, lo que impide
su uso como control preventivo en el pipeline \acs{CICD}.

La Tabla~\ref{tab:comparativa} resume el posicionamiento de \textit{kubescan} frente
a las principales herramientas del estado del arte.

\begin{table}[htbp]
\centering\small
\caption{Comparativa de herramientas de seguridad para Kubernetes}
\label{tab:comparativa}
\begin{tabular}{@{}lcccc@{}}
\toprule
Herramienta & \shortstack{Análisis\\estático} & \shortstack{Requiere\\clúster} & Grafos & \shortstack{\ac{ML}\\supervisado} \\
\midrule
Checkov \citep{Checkov2021} & \checkmark & No & Parcial & No \\
kube-linter \citep{KubeLinter2021} & \checkmark & No & No & No \\
Kubescape \citep{Kubescape2021} & \checkmark & Opcional & No & No \\
KubeHound \citep{KubeHound2023} & No & Sí & \checkmark & No \\
IceKube \citep{IceKube2023} & No & Sí & \checkmark & No \\
\midrule
\textbf{kubescan} & \checkmark & \textbf{No} & \checkmark & \checkmark \\
\bottomrule
\end{tabular}
\fuentepropia
\end{table}

El sistema \textit{kubescan} presentado en este trabajo responde a las tres carencias
identificadas: adopta una arquitectura de grafo heterogéneo con tipos de arista
semánticamente significativos para el dominio \textit{Kubernetes}, combina el análisis
estático (Capa~1) con el razonamiento estructural (Capa~2) en un \textit{ensemble}
optimizado (Capa~3), y contribuye un conjunto de datos de 599~grafos de clúster
originales con etiquetas de cadena de ataque derivadas de reglas de seguridad
explícitas y auditables. Al operar sobre manifiestos \acs{YAML} sin requerir acceso
al clúster, \textit{kubescan} se posiciona como un control preventivo integrable en
pipelines \acs{CICD}, complementario a —y no sustituto de— las herramientas de
análisis de clúster activo como KubeHound.
